\documentclass[../main.tex]{subfiles}
\begin{document}
%==========================================TIÊU ĐỀ TẬP========================================
\begin{table}[h]
  \begin{adjustwidth}{-1cm}{}
    \begin{tabular}{>{\centering\arraybackslash}p{6cm}|>{\raggedright\arraybackslash}p{12cm}}
      \multirow{5}{6cm}
      % Cột bên trái: ÉPISODE và số tập
      {\\[-40pt]
        \centering
        {\fontsize{20pt}{20pt}\selectfont{\outlinetext[1]{eptype}{white}{\flaregothic PERSONNALISÉ}}\\[12pt]
          {\fontsize{36pt}{36pt}\selectfont{\outlinetext[1]{epnum}{white}{\burbank CHÍN}}}\\[6pt]     % use for custom episodes
          {\fontsize{20pt}{20pt}\selectfont{\outlinetext[1]{epnum}{white}{\cafeteria (C-9)}}}\color{black}}
        \\[18pt]
        % {\fontsize{14pt}{14pt}\selectfont{\outlinetext[1]{epattr}{white}{\flaregothic [I]}}}\\
        {\fontsize{18pt}{18pt}\selectfont{\outlinetext[1]{epattr}{white}{\sen Mini-histoire}}}\\
        \color{black}
      }
       &
      % Dòng thứ nhất: tên tập tiếng Nhật
      {\fotskip\fontsize{18pt}{18pt}\selectfont \ruby{越}{えつ}\ruby{南}{なん}\ruby{料}{りょう}\ruby{理}{り}} \\[6pt]
      % Dòng thứ hai: tên tập tiếng Anh
       & {\cafeteria\fontsize{18pt}{18pt}\selectfont Mixed Food Course}                           \\[4pt]
      % Dòng thứ ba: tên tập tiếng Việt
       & {\vnmsans\fontsize{18pt}{18pt}\selectfont Những món ăn độc nhất vô nhị xứ Etsunan}       \\[8pt]
      % Dòng thứ tư: Nhân vật xuất hiện
       & {\caviar{\fontsize{14pt}{14pt}\selectfont Nhân vật: Theo từng câu chuyện}}               \\
      % Dòng thứ năm (nếu có): Nhân vật chỉ xuất hiện trong Omake
      %  & {\abv Truyện phụ:} \emp{okayu} 
      \\[8pt]
      % Dòng cuối cùng: Thời gian diễn ra của tập (thường sẽ là ngày phát hành)
      %  & {\sen\fontsize{16pt}{16pt}\selectfont Mốc thời gian: DD/MM/YYYY}                                     \\[18pt]
    \end{tabular}\\[8pt]
  \end{adjustwidth}
\end{table}
%=========================================TRUYỆN CHÍNH========================================
\color{black}

\txtbx[2]{
  {\color{vol} \uline{Tóm tắt:}} Tuyển tập 9 mẩu chuyện ngắn về ẩm thực dân dã xứ Etsunan do Kuon và các em gái nhà Hikawa giới thiệu. Từ những món ăn vặt tuổi thơ như bánh mì chấm sữa đặc, xoài xanh chấm muối ớt, bánh tráng nướng, nước mía, sa tế tôm, chè hạt sen, cho đến những đặc sản ``kinh dị'' hơn như mắm tôm, trứng vịt lộn và chả rươi. Mỗi món đều khiến các cô gái \hll\ đi từ hoang mang, sợ hãi đến bất ngờ và mê mẩn. Flare, Towa, Aki, Ayame liên tục bị ``dắt mũi'' qua từng trải nghiệm ẩm thực, trong khi Suzuno và Ryuuhou hào hứng phụ họa cùng Kuon. Kết thúc, tất cả đều công nhận rằng ẩm thực quê hương cậu Tổng tuy nhìn đơn giản nhưng lại có sức hút khó cưỡng.
}

Nền ẩm thực Etsunan từ lâu đã nổi tiếng khắp thế giới với sự phong phú, tinh tế và cân bằng hương vị. Người ta thường nhắc đến những bát phở bò thơm phức nghi ngút khói, những ổ bánh mì kẹp thịt giòn rụm đậm đà, hay bún chả đậm đà hương vị truyền thống.

Tuy nhiên, thay vì mang đến văn phòng COVER những món ăn tráng lệ hay đã quá quen thuộc trên các tạp chí du lịch, Kuon lại có một ý tưởng kỳ lạ: giới thiệu những ``huyền thoại'' bình dị, độc đáo và đầy ký ướm tuổi thơ của người dân quê nhà, đôi lúc còn khiến thực khách phải ``mắt chữ A, mồm chữ O''.

Đó là những sự kết hợp tưởng chừng như ngẫu hứng, giản đơn đến mức tối giản, nhưng lại ẩn chứa nét văn hóa đời thường vô cùng sinh động.Và lần này, với sự hỗ trợ từ các em gái nhà Hikawa, cậu Tổng lại tiếp tục mang đến cho các cô gái \hll\ những trải nghiệm ẩm thực có một không hai.

\vspace{-8pt}

\subsection*{\phantomsection\label{ep-C9.I}}
\addcontentsline{toc}{subsection}{{\sen{-Histoire 1-}} {\caviar Bánh mì chấm sữa đặc}}
\stpt{-Histoire 1-}
\stcap{Bánh mì chấm sữa đặc}
\vspace{-4pt}
\begin{center}
  {\caviar\fontsize{12pt}{12pt}\selectfont Nhân vật: \emp{kuon2} \emp{suzuno2} \emp{aki} \emp{flare1} \emp{towa}}
\end{center}

Một ngày như mọi ngày tại văn phòng \hll, không khí mát rượi từ chiếc điều hòa cũng không ngăn được sự tò mò đang dâng lên trong góc phòng nghỉ.

Kuon ngồi vắt chân trên ghế tựa, tay cầm tách trà ấm, cất giọng đố đầy vẻ bí hiểm: ``Mấy đứa này, đố biết người đàn ông nào trên thế gian này lại có sữa?''

Ba cô gái gồm Flare, Towa và Aki ngơ ngác nhìn nhau trước câu hỏi có phần vô lý này. Suzuno đang đứng bên kệ bếp gần đó dọn dẹp ly tách cũng phải ngẩng đầu lên, chớp chớp đôi mắt dị sắc lấp lánh sự tò mò.

Nàng Dị giới đưa tay lên cằm, lập tức nghiêng đầu đáp: ``\dots Không có ai cả chứ sao anh? Làm gì có chuyện kỳ quặc đó.''

Nàng Ác ma gật đầu lia lịa, đuôi quỷ sau lưng khẽ vẫy: ``Đúng đó! Em chưa bao giờ nghe nói hay thấy người đàn ông nào lại có sữa cả.''

Nàng quỷ Thế hệ thứ Tư thậm chí còn liên tưởng đến hình ảnh mấy con quỷ Succubus trong sách cổ, má hơi đỏ lên rồi xua tay: ``Em nghĩ chắc chỉ có phụ nữ mới có thể cho ra sữa thôi chứ!''

Aki mỉm cười gật gù: ``Giống như chị em mình đúng không?''

``Đấy, kiểu kiểu thế đó!'' Towa tán thành.

Khác với hai người kia, Flare khẽ vuốt cằm, trầm ngâm một lúc rồi đưa ra một thông tin mang tính nhân học đầy bất ngờ: ``Chị nhớ không nhầm thì tộc Aka ở châu Phi có những người đàn ông có thể cho con bú được đấy.''

Aki và Towa quay sang nhìn Flare với bản mặt ngạc nhiên tột độ: ``Thật luôn!? Có cả chuyện kỳ diệu đó sao, Flare-chan/-senpai?''

Kuon phì cười, vội lắc đầu đính chính: ``Flare-chan nhớ gần đúng rồi đấy, nhưng thực ra đàn ông tộc Aka chỉ đưa ti ra cho em bé bú dỗ để dỗ chúng nín khóc lúc người mẹ đi vắng thôi, chứ cơ thể nam giới ở đâu cũng không thực sự tiết ra sữa được đâu.''

Cậu khẽ nhấp một ngụm trà rồi tiếp tục: ``Và đó cũng không phải là câu trả lời anh đang muốn hướng tới.''

Cả ba cô gái trơ mặt ra. Towa thắc mắc: ``Thế rốt cuộc câu trả lời là gì vậy, Kuon-senpai?''

Không nói không rằng, cậu Tổng thò tay vào túi xách đặt cạnh bàn, lấy ra một lon sữa đặc bằng thiếc đặt lên bàn. Trên vỏ lon in hình một ông lão râu tóc bạc phơ trông giống như một vị thần tiên, đi kèm dòng chữ nổi bật: ``Trường Thọ''.

``Sữa đặc `Trường Thọ' ư!?'' ba cô gái đồng thanh reo lên.

Kuon vỗ bốp một cái lên lon sữa: ``Chính nó. Mấy đứa đã từng thấy lon này bao giờ chưa?''

Aki tò mò ghé sát lại quan sát: ``\dots Chưa ạ. Nhìn cũng chỉ giống mấy lon sữa đặc trong siêu thị thôi mà.''

Towa lắc đầu: ``Em cũng chưa thấy bao giờ.''

Thế nhưng, Flare lại trố mắt, reo lên đầy kinh ngạc: ``\dots Em biết cái này!''

Aki và Towa lập tức chĩa ánh mắt khó hiểu về phía cô. Towa ngơ ngác: ``Chị biết lon này thật sao Flare-senpai?''

Nàng Yêu tinh vàng gật đầu chắc nịch: ``Chị đã từng thấy nó rồi\dots\ Sữa đặc này là để chấm với bánh mì đúng không!?''

Hai cô nàng tròn mắt kinh ngạc: ``Thật á!?''

Ngay lúc đó, Suzuno mỉm cười bước lại gần, dịu dàng giải thích thêm cho mọi người: ``Flare-san giỏi thật đấy ạ! Ở quê em, đây chính là món quà sáng giá rẻ và cũng là đồ ăn vặt `huyền thoại' của học sinh, sinh viên đấy ạ. Chỉ cần một ổ bánh mì nóng giòn và một ít sữa đặc ngọt béo là đã đủ năng lượng cho cả buổi sáng rồi. Onii-sama và Onee-sama hồi trước toàn mua món này về cho em với Ryuu ăn mỗi khi thèm ngọt đấy.''

Kuon gật gù hài lòng: ``Suzuno nói đúng đấy. Hồi ngày xưa ở nhà, bọn anh hay ăn món này lắm. Nhanh, rẻ mà lại ngon cực kỳ.''

Flare vui vẻ đập tay: ``Giờ chỉ còn thiếu mỗi bánh mì nữa là đủ bộ luôn rồirồi\dots!''

Cậu Tổng chỉ tay về phía bếp: ``Anh có treo hai ổ bánh mì dài ở trên kệ bếp đấy.''

``Thế á!? Để em lấy cho!'' nàng Yêu tinh hào hứng đứng phắt dậy đi thẳng vào bếp, để lại Aki và Towa dõi theo bóng lưng cô với sự ngỡ ngàng.

``Chị ấy trông hào hứng ghê\dots'' nàng Ác ma thì thầm.

Chỉ một lúc sau, Flare quay lại với hai ổ bánh mì baguette giòn rụm trên tay. Towa nhìn ổ bánh mì rồi lại nhìn lon sữa, hơi lo lắng: ``Kuon-senpai\dots\ Anh có chắc là kết hợp thế này ăn được không vậy?''

Kuon thong thả lấy đồ khui mở nắp lon sữa: ``Cứ thử đi rồi biết, không ngon anh làm con ngươi.''

Aki bê đĩa lại gần: ``Chị nghĩ là ngon đấy, thử xem sao!''

Cậu Tổng xé ra hai miếng bánh mì lớn rồi đưa cho từng người. Aki lập tức cầm miếng bánh mì nóng giòn, chấm đẫm vào lòng lon sữa đặc sánh mịn vàng ươm.

``E-Em ăn nha!''

Nàng Dị giới kêu lên một tiếng rồi cắn trọn miếng bánh mì vào miệng. Ngay lập tức, một luồng ánh sáng thánh thiện tỏa ra xung quanh người cô như một thiên thần đang giang cánh. Aki chắp hai tay lên má, mắt nhắm tịt vì sung sướng: ``Ngon quớ đi\~{}\~{}!''

Nàng Ác ma giật mình trước phản ứng mãnh liệt của tiền bối. Không thể chịu nổi sự kích thích, cô nàng cũng vội vàng xé bánh, chấm một lượng sữa thật lớn: ``E-Em cũng muốn thử nữa!''

Cô cho nguyên miếng bánh vào miệng. Rập khuôn phản ứng của Aki, một vầng sáng thiên đường bùng nổ quanh người nàng quỷ: ``Tuyệt cú mèo\~{}\~{}!!''

Flare đứng bên cạnh gật gù hài lòng: ``Đúng chứ? Cái cảm giác vỏ bánh mì giòn rụm kết hợp với vị sữa đặc béo ngậy, ngọt ngào bùng nổ trong khoang miệng\dots\ nó khó tả lắm đúng không? Vì chị cũng từng trải qua cảm giác đó rồi mà\dots''

Chưa kịp để cô nàng trải lòng hết cảm xúc, Kuon đã bật cười lớn: ``Flare-chan, em mê mẩn tâm sự cũng được thôi, nhưng nhìn lại kìa, mấy đứa kia chén sạch số bánh mì rồi đấy.''

Nàng Yêu tinh vàng khựng lại. Cô từ từ quay mặt sang bàn ăn, và đập vào mắt là cảnh tượng hai ổ bánh mì dài ngoằng đã biến mất không còn một vết tích từ lúc nào. Bên cạnh đó, không chỉ Aki và Towa mà ngay cả Suzuno - với tâm hồn ăn uống vô đáy bẩm sinh - cũng đang cầm khăn lau mép, nét mặt mãn nguyện vô cùng.

``\textbf{T-Tại sao chứ!?}'' Flare thảng thốt kêu lên, hai tay ôm đầu đầy thất vọng, ``Chị còn chưa kịp ăn miếng nào mà!''

Đáp lại sự uất hận của cô nàng Thế hệ thứ Ba là những nụ cười trừ từ Aki, Towa và Suzuno. Cả ba miệng vẫn còn dính chút sữa đặc và vụn bánh mì.

``Hì hì, em xin lỗi nhé, Flare-san\dots\ Tại bánh mì giòn quá mà sữa đặc mà Onii-sama mang tới ngon thật sự luôn ấy!'' Suzuno lễ phép gãi đầu xua tay.

``Đúng đó, ngon không chịu nổi luôn!'' Towa tiếp lời.

May mắn cho Flare, Kuon thở dài lắc đầu rồi mở tủ lạnh lấy ra thêm hai ổ bánh mì dự phòng mà cậu mua về ăn dần. Quả nhiên, đối với món ăn gây nghiện này, kể cả có đưa cho họ mười ổ bánh mì thì bụng họ vẫn cảm thấy thèm!

\newpage

\subsection*{\phantomsection\label{ep-C9.II}}
\addcontentsline{toc}{subsection}{{\sen{-Histoire 2-}} {\caviar Xoài xanh chấm muối ớt}}
\stpt{-Histoire 2-}
\stcap{Xoài xanh chấm muối ớt}
\vspace{-4pt}
\begin{center}
  {\caviar\fontsize{12pt}{12pt}\selectfont Nhân vật: \emp{kuon2} \emp{ryuuhou2} \emp{ayame} \emp{flare1} \emp{towa}}
\end{center}

Vào một buổi sáng khác tại phòng làm việc, Kuon xách theo một túi nilon chứa ba quả xoài xanh ươm, to tròn. Cậu thong thả đặt đĩa xoài lên bàn tròn giữa phòng rồi quay lưng đi vào bếp tìm dụng cụ.

Đúng lúc này, Ayame và Towa cùng nhau bước vào văn phòng.

``Chào buổi sáng-yo!'' nàng Quỷ sừng reo lên đầy năng lượng. Nhưng ngay lập tức, ánh mắt cô bị thu hút bởi ba quả quả lạ màu xanh lục nằm trên đĩa.

Ayame tiến lại gần, nghiêng đầu: ``À-rế? Cái gì đây nhỉ?''

Towa nheo mắt nhìn kĩ: ``Trông chúng giống như mấy quả xoài thì phải.''

Nàng Oni ngơ ngác: ``Xoài\dots? Nó là loại quả gì vậy Towa-chan?''

Nàng Ác quỷ lục lại ký ức rồi giải thích: ``Em nghĩ đó là một loại quả có vỏ màu vàng, ăn khá ngọt\dots\ Chắc thế? Người ta hay dùng nó để làm sinh tố hoặc đồ uống ấy.''

Ojou-san cúi xuống nhìn kỹ ba quả xoài cứng ngắc màu xanh lá cây: ``Nhưng mà đây là xoài xanh lè mà?''

Ryuuhou vốn đang ngồi vắt chân trên chiếc sofa góc phòng ôm điện thoại liền ngẩng đầu lên, chêm vào một câu bằng giọng điệu vô cùng am hiểu: ``Bởi vì đó là xoài chưa chín đấy hai chị ạ! Nhưng ở quê em, xoài xanh thế này mới là đỉnh cao ăn vặt đấy!''

Cùng lúc đó, Flare từ phòng trà bước ra: ``Đúng như Ryuu-chan nói đấy, nó là xoài chưa chín!''

Ayame giật mình quay lại: ``Ái chà, Flare! Cả Ryuu-chan nữa!''

Towa thở tháo: ``Flare-senpai\dots! Làm em giật cả mình!''

Nàng Oni chỉ vào quả xoài trong túi nilon kia: ``Em bảo xoài này chưa chín\dots\ Vậy nó có ăn được trực tiếp không?''

Flare bước lại gần: ``Về lý thuyết thì nó hoàn toàn ăn được-''

Chưa kịp để cô nói hết câu, Ayame với tính cách hấp tấp vốn có đã chộp ngay lấy một quả xoài, đưa lên miệng cắn một miếng thật to. ``Vậy mời mọi người nha!''

Towa tá hỏa giơ tay cản: ``Ơ, khoan đã Ayame-senpai!''

*RẮC!*

Vừa cắn rứt một miếng xoài sống còn nguyên vỏ, nét mặt Ayame lập tức biến sắc. Vị chua chát xộc thẳng lên não khiến cả người cô run rẩy rồi sụp đổ hoàn toàn xuống sàn nhà, hai tay ôm lấy má: ``C-Chua quáaa\dots\ Răng chị tê rần luôn rồi nè!''

Flare và Towa đổ mồ hôi hột trước cảnh tượng đó. Nàng Yêu tinh vàng thở dài: ``Em còn chưa kịp nói hết mà\dots\ Xoài sống chua lắm, khó ăn trực tiếp thế này lắm chị ơi.''

Nàng Ác ma bất lực đỡ trán: ``Với cả quả xoài còn chưa rửa, chưa gọt vỏ nữa chứ\dots''

Ryuuhou ngồi trên sofa không tha cơ hội trêu chọc, cô bé phì cười: ``Ayame-san bá đạo thật đấy, cắn nguyên quả chưa gọt vỏ luôn! Bái phục Ojou-san luôn rồi!''

``Thế nên món này nó mới phải có `vũ khí đi kèm' để chấm vào đấy chứ.'' Từ đằng sau ba người, Kuon bước ra với một con dao gọt hoa quả, một chiếc đĩa sứ sạch và một bát muối ớt đỏ rực đầy hấp dẫn.

Towa quay lại, nhìn con dao mà cậu đang cầm trên tay, kêu lên: ``Kuon-senpai! Lại là món kỳ lạ gì của anh nữa đây!?''

Flare mỉm cười hào hứng: ``Kuon-senpai, anh lại mang món gì đó hay ho từ quê nhà sang sao?''

Kuon đặt đồ lên bàn, lắc đầu nhìn nàng Ác ma: ``Tất nhiên rồi. Mà Towa-chan này, quỷ gì mà yếu bóng vía thế, thấy dao kéo đã giật nẩy mình lên rồi.''

``E-Em không có!'' cô nàng đỏ mặt cãi lại, vội vàng cúi xuống dìu Ayame đang nhăn nhó lên chiếc ghế bành.

\ct

Sau khi Kuon nhanh nhẹn gọt sạch vỏ và thái ba quả xoài xanh thành những miếng nhỏ vừa ăn, đĩa xoài xanh giòn sần sật đã sẵn sàng.

Cả ba người Flare, Towa, Ayame cùng cất tiếng hỏi: ``Xoài xanh chấm muối ớt?''

Cậu Tổng gật đầu: ``Ừ. Lại là một món ăn vặt dã ngoại thần thánh ở quê anh.''

Ryuuhou liền nhanh nhảu bước lại gần đĩa xoài, vừa dùng tăm xiên một miếng vừa hào hứng giải thích tiếp lời anh trai: ``Để em giải thích cho! Xoài xanh chua lè ăn một mình thì khó nuốt thật, nhưng khi chấm với muối ớt mặn mặn, cay cay và ngọt nhẹ thì vị chua sẽ bị `khống chế' hoàn toàn. Bọn học sinh bên quê em đứa nào cũng mê món này, cứ ra cổng trường là phải làm một bịch ăn rôm rả với nhau mới chịu được!''

Nói rồi, cô chấm nhẹ miếng xoài vào bát muối ớt rồi cho vào miệng cắn *RỘP* một tiếng đầy khoái chí: ``Đấy, các chị nghe tiếng giòn rụm này đi, bùng nổ hương vị luôn!''

Kuon cũng thong thả lấy một miếng ăn thử: ``Ryuuhou nói chuẩn đấy. Món này ăn lúc rảnh rỗi là hết bài.''

Flare nhìn hai anh em ăn ngon lành liền gật gù: ``Đồ ăn chua đi với đồ mặn cay đúng là một sự kết hợp rất khoa học và hợp lý.'' Nàng Yêu tinh vàng thoải mái nhón một miếng ăn thử.

Bất chợt, Ayame đứng bên cạnh nhìn đĩa xoài mà khóe miệng chảy nước miếng ừng ực. Towa giật thót mình lo lắng: ``A-Ayame-senpai? Chị sao thế?''

Nàng Quỷ sừng nuốt nước bọt, mắt dính chặt vào đĩa xoài: ``Chị\dots\ chị cũng muốn thử lại một miếng nữa\dots''

Nàng Ác ma hoảng hốt cản lại: ``Chị mới vừa nhè ra nhăn nhó xong đấy thôi!''

Ayame quay sang cô, ánh mắt long lanh đầy thèm thuồng: ``Nhưng mà Towa-chan à\~{} Em nghe tiếng `rộp rộp' khi anh Kuon với Ryuu-chan nhai có thấy kích thích không?''

Towa ngập ngừng, nhìn miếng xoài giòn ươm trên đĩa: ``Đ-Đúng là nghe giòn thật\dots\ Nhưng chưa chắc nó đã ngon đâu! Vừa nãy chị chua đến quéo cả lưỡi rồi còn gì!''

``Cái này cũng giống như nhìn người ta ăn chanh thôi mà, tự nhiên thèm lắm luôn ấy!''

``Chị đang nói cái gì kỳ cục vậy chứ!?''

Thấy nàng Ác quỷ vẫn còn lưỡng lự, Flare chỉ tay vào đĩa xoài, ôn tồn khuyên: ``Hai người cứ ăn thử theo cách Kuon-senpai với Ryuu-chan chỉ đi. Nó không đáng sợ như hai người nghĩ ban đầu đâu.''

Nhận được sự tiếp sức, Ayame ngay lập tức nhón lấy một miếng xoài thật to. Thế nhưng vì quá nôn nóng, cô ấn mạnh miếng xoài xuống bát muối ớt làm nó dính một lớp muối dày cộp!

Cậu Tổng vội vàng can: ``Ấy! Chấm nhẹ tay thôi em!''

Nhưng đã quá trễ, Ayame đã đưa nguyên miếng xoài đầy muối vào miệng cắn *RẮC*!

Và\dots\ không ngoài dự đoán, nàng Quỷ sừng lại một lần nữa nhè miếng xoài ra, mặt mày biến dạng vì mặn chát. Towa ngơ ngác, còn Kuon, Flare và Ryuuhou thì đồng loạt đổ mồ hôi hột.

``M-Mặn quáaaa! Cay xè lưỡi luôn rồi-yo!'' nàng Oni lè lưỡi khóc dở mếu dở.

Towa cào đầu bứt tóc: ``A-Ayame-senpai!! Đã bảo chậm chậm thôi mà!''

Ryuuhou thở dài thè lưỡi: ``Trời ạ, đã dặn là chấm nhẹ cái đầu miếng xoài thôi mà\dots''

Flare lắc đầu mỉm cười: ``Chị đừng có vội vã quá thế chứ. Towa-chan, em thử làm mẫu cho chị ấy xem đi?''

Nàng Ác quỷ ngập ngừng tiến lại gần đĩa xoài. Cô chọn một miếng xoài nhỏ, cẩn thận quệt nhẹ một chút xíu muối ớt ở góc theo đúng lời dặn.

``\textit{L-Liệu hai anh em họ có đang chơi khăm mình không đây\dots}'' Towa thầm nghĩ trong đầu, nhắm mắt nhắm mũi cắn một miếng.

*RỘP!*

Tiếng nhai giòn tan vang lên. Ngay tức khắc, đôi mắt xanh lá mạ của Towa mở to kinh ngạc. Vị chua thanh nguyên bản của xoài hòa quyện hoàn hảo với vị mặn ngọt và cay nồng của muối ớt tạo nên một cơn bùng nổ hương vị cực kỳ kích thích.

``Cũng\dots\ cũng không đến nỗi nào nhỉ! Không, phải nói là ngon ấy chứ!'' nàng Ác ma trầm ồ khen ngợi, rồi cô liên tiếp chấm thêm miếng thứ hai, thứ ba. Chẳng mấy chốc, miếng xoài ban đầu đã nằm gọn trong bụng cô.

Ayame thấy hậu bối ăn ngon lành liền xóc lại tinh thần. Cô cẩn thận gắp miếng khác, bắt chước từng thao tác nhẹ nhàng của Towa: quệt nhẹ một chút muối ớt.

*RỘP!*

Một tiếng cắn nữa vang lên. Trong khoang miệng của cô, vị chua mặn cay hòa quyện bùng nổ. Đôi mắt nàng quỷ sừng sáng rực lên như đèn pha, cô reo lên tấm tắc: ``Ngon thế! Giòn sần sật mà cay cay đã quá đi-yo!''

Ryuuhou khoanh tay mỉm cười đắc thắng: ``Em đã bảo rồi mà! Đồ ăn vặt quê em tuy nhìn đơn giản nhưng ăn cuốn lắm!''

Flare nhẹ nhàng nhón thêm miếng nữa, mỉm cười: ``Thấy chưa? Nó đâu có đáng sợ như mọi người nghĩ đâu.''

Kuon nhìn đĩa xoài bắt đầu vơi đi nhanh chóng, mỉm cười xoa đầu em gái rồi quay sang các holomem: ``Trong túi anh còn nhiều xoài lắm, mấy đứa cứ ăn thoải mái đi!''

Và thế là, thêm một món ăn vặt bình dị, dân dã nhưng đầy ma lực nữa từ quê hương cậu Tổng đã chính thức ghi tên vào danh sách món khoái khẩu của các thành viên \hll.

\newpage

\subsection*{\phantomsection\label{ep-C9.III}}
\addcontentsline{toc}{subsection}{{\sen{-Histoire 3-}} {\caviar Bánh tráng nướng}}
\stpt{-Histoire 3-}
\stcap{Bánh tráng nướng}
\vspace{-4pt}
\begin{center}
  {\caviar\fontsize{12pt}{12pt}\selectfont Nhân vật: \emp{kuon2} \emp{ryuuhou2} \emp{aki} \emp{ayame}}
\end{center}

Vào một buổi chiều khác, Kuon bước vào phòng nghỉ của văn phòng với một hộp đồ ăn bọc trong lớp giấy bạc cẩn thận để giữ nhiệt. Đi ngay bên cạnh cậu là Ryuuhou đang lanh chanh xách theo một túi đĩa nhựa và khăn giấy.

Mừa mở cửa phòng, một mùi thơm nức mũi của hành phi, bơ, thịt bằm và sốt nướng lập tức lan tỏa khắp không gian, thu hút sự chú ý của Aki và Ayame đang ngồi ở bàn trà.

Nàng Dị giới hít hà một hơi, đôi mắt sáng rực reo lên: ``Mùi này\dots\ Quen quá! Pizza đúng không anh!?''

Nàng Quỷ sừng vừa nghe thấy từ ``pizza'' cũng lập tức bật dậy, hai tay giơ lên trời hô vang như hô khẩu hiệu: ``Pizza! Pizza! Pizza-yo!''

Kuon phì cười, lắc đầu đáp: ``Không hẳn là pizza đâu. Nhưng mà cũng gần giống như thế.''

Ryuuhou đặt túi đĩa xuống bàn, khoanh tay mỉm cười trêu chọc: ``Gần đúng thôi hai chị ơi! Đây là `Pizza Etsunan' đó!''

Hai cô gái giương mắt nhìn nhau đầy khó hiểu. Không để họ chờ lâu, Cậu Tổng nhẹ nhàng mở chiếc hộp giấy bạc ra. Bên trong là ba chiếc bánh tráng nướng nóng hổi, giòn rụm với bán kính cỡ một gang tay. Mỗi chiếc mang một loại nhân khác nhau: một chiếc nhân xúc xích thái mỏng, một chiếc phủ đầy trứng cút vàng ươm, và chiếc còn lại ngập tràn thịt bằm thơm phức. Bề mặt cả ba chiếc bánh đều được rưới những đường sốt tương cà, tương ớt và sốt mayonnaise béo ngậy đẹp mắt.

Aki ghé sát lại quan sát: ``Sao trông chúng giống pizza thế nhỉ?''

Ayame gật gù: ``Mùi cũng giông giống nữa\dots\ Rốt cuộc nó là món gì vậy, Kuon-senpai?''

Kuon mỉm cười giới thiệu: ``Bánh tráng nướng. Đây là món ăn vặt đặc sản cực kỳ nổi tiếng của Tagaku (Đà Lạt) đấy.''

Hai người con gái lại ngơ ngác nhìn nhau. Aki chớp mắt thắc mắc: ``Tagaku? Nó là ở đâu vậy anh?''

Ryuuhou ôn tồn giải thích: ``Đó là một thành phố nằm ở vùng cao nguyên phía Nam Etsunan. Nhờ địa hình trên cao nên không khí ở đó quanh năm mát lạnh, dễ chịu lắm. Em từng đi chơi với mẹ ở đó, nên em biết!''

Aki tò mò hỏi tiếp: ``Trên đó có tuyết rơi không?''

Ayame lấp lánh đôi mắt: ``Nghe thích quá! Em cũng muốn đến đó chơi-yo!''

Kuon đổ mồ hôi hột cười trừ: ``Ở đó không lạnh đến mức có tuyết đâu. Nhưng bù lại, thời tiết se lạnh đó mà thưởng thức ẩm thực thì tuyệt vời lắm.''

Em gái cậu đứng bên cạnh tiếp lờii, giọng điệu đầy tự hào: ``Đúng đó hai chị! Ở Tagaku, mỗi khi chiều tối trời bắt đầu trở lạnh, người ta thường ngồi quanh các bếp than hồng ngoài phố. Vừa nhìn người bán hàng nướng chiếc bánh tráng giòn rụm vừa thổi ăn ngay tại chỗ thì chỉ có `hết nước chấm'!''

Nàng Dị giới chỉ vào ba chiếc bánh trên bàn: ``Và đây là một trong những món đó sao?''

Kuon gật đầu. Ayame thắc mắc: ``Nhưng sao anh lại gọi nó là `Pizza Etsunan' chứ?''

Kuon đáp: ``Vì hình dáng tròn trịa, có phần đế và phủ nhiều loại nhân lên trên rưới thêm sốt khá giống pizza Ý, nên khách du lịch gọi vui như thế. Nhưng phần đế này làm từ bánh tráng mỏng nướng trên than hồng nên vị của nó giòn tan và rất khác biệt đấy.''

Thấy Aki và Ayame nhìn ba chiếc bánh nóng hổi có phần do dự, cậu Tổng xua tay: ``Nếu không tin thì hai đứa cứ ăn thử đi. Đằng nào một mình anh cũng chẳng ăn hết nổi cả ba cái này đâu.''

Nàng Oni hào hứng giơ tay ngay tắp lừ: ``Vậy em xí cái bánh nhân trứng cút nha!''

Aki cũng nhanh tay chọn: ``Chị lấy cái nhân thịt bằm vậy!''

Ryuuhou nhanh nhẹn chia đĩa cho từng người. Kuon hỏi: ``Có cần anh lấy kéo cắt cho không?''

Nàng Quỷ sừng lắc đầu bạo dạn. Cô rút chiếc dao nhỏ đi kèm bên hông, nhanh như chớp cắt chiếc bánh ra làm tư, trong khi Aki cũng khéo léo dùng năng lượng ma thuật dị giới để chia chiếc bánh của mình thành bốn phần bằng nhau.

Mỗi người cầm một miếng bánh tráng nướng nóng hổi lên. Cả hai đồng thanh lễ phép: ``M-Mời cả nhà ăn cơm! / Mời cả nhà ăn cơm-yo!!''

*RẮC!*

Một tiếng cắn giòn tan vang lên. Ngay khi miếng bánh chạm vào lưỡi, đôi mắt của cả Aki và Ayame đều mở to kinh ngạc. Sự kết hợp giữa vị giòn rụm của bánh tráng nướng bơ, vị béo ngậy của sốt mayonnaise, cay nhẹ của tương ớt và đậm đà của phần nhân khiến khoang miệng họ như bùng nổ.

``Ngon thế này!!'' cả hai thốt lên đầy trầm ồ.

Nàng Dị giới chắp tay lên má mãn nguyện: ``Độ giòn của vỏ bánh mì tráng này thật kỳ diệu\dots''

Nàng Oni gật gù liên tục: ``Quyện cùng nước sốt béo béo cay cay rưới ở trên nữa chứ!''

``Nhân thịt bằm đậm đà, mềm mại vừa vặn cực kỳ!''

``Nhân trứng cút bùi bùi cũng bùng nổ lắm nè-yo!''

Cả hai cùng đồng thanh: ``Lại còn có cả mùi thơm của hành lá nướng nữa chứ!''

Chẳng mấy chốc, hai chiếc bánh tráng nướng đã biến mất hoàn toàn vào bụng hai cô gái. Thế nhưng, vừa nuốt xong miếng cuối cùng, họ lập tức quay sang\dots\ tranh cãi chí chóe xem loại nhân nào đỉnh hơn.

Aki khoanh tay khẳng định: ``Chắc chắn là nhân thịt bằm ngon hơn rồi!''

Ayame bĩu môi cãi lại: ``Không đúng! Trứng cút béo ngậy mới là chân lý chứ-yo!''

``\textbf{Thịt bằm!}''

``\textbf{Trứng cút!}''

Thấy cuộc chiến chuẩn bị leo thang, cậu Tổng vội vàng giơ hai tay can thiệp: ``Thôi thôi, dừng lại nào hai đứa! Nhân nào cũng có cái ngon riêng của nó mà.''

Ryuuhou đứng bên cạnh phì cười, tủm tỉm mách nước: ``Đúng đó, hai chị cứ mải cãi nhau đi, để em với Onii xử nốt cái bánh nhân xúc xích còn lại là vừa đẹp!''

Nghe thấy thế, Aki và Ayame giật mình quay lại nhìn chiếc bánh cuối cùng trên đĩa. Nàng Quỷ sừng nuốt nước bọt hỏi: ``Kuon-senpai\dots\ Ăn món này giống pizza thật đấy, vị lại còn rất mới lạ nữa\dots\ Hôm nào đó anh làm thêm cho bọn em ăn nữa nha?''

Nàng Dị giới chắp hai tay trước ngực, mắt long lanh: ``Đúng đó anh, làm cho bọn em ăn nữa đi ạ!''

Kuon đổ mồ hôi hột nhìn hai cô gái vừa ăn sạch hai cái bánh mà vẫn còn thèm thuồng: ``Hai đứa đúng là dễ bị mua chuộc thật đấy\dots'' Nhưng rồi cậu cũng mỉm cười xoa đầu cả hai, ``Được rồi, hôm nào rảnh anh và Ryuuhou sẽ làm một chầu bánh tráng nướng hoành tráng cho cả văn phòng.''

Nghe đến đây, hai người vui sướng reo hò ầm ĩ. Còn về phần cậu Tổng và em gái cậu\dots\ họ lập tức nhanh tay chộp lấy miếng bánh nhân xúc xích cuối cùng cho vào miệng, bởi họ thừa biết với độ phàm ăn của hai cô nàng này, chậm chân một giây thôi là đến cả họ cũng chẳng còn mẩu vụn nào!

\newpage

\subsection*{\phantomsection\label{ep-C9.IV}}
\addcontentsline{toc}{subsection}{{\sen{-Histoire 4-}} {\caviar Nước mía}}
\stpt{-Histoire 4-}
\stcap{Nước mía}
\vspace{-4pt}
\begin{center}
  {\caviar\fontsize{12pt}{12pt}\selectfont Nhân vật: \emp{kuon2} \emp{kaigou2} \emp{naomi2} \emp{aki} \emp{ayame} \emp{flare1} \emp{towa}}
\end{center}

Một ngày nắng nóng bất thường vào giữa mùa xuân. Thời tiết oi bức đột ngột khiến không khí trong văn phòng \hll\ trở nên ngột ngạt.

Trong phòng nghỉ, Towa và Flare đang nằm sõng soài trên chiếc ghế sofa dài, toàn thân uể oải không chút sức sống vì cái nóng.

Towa than thở, chiếc đuôi quỷ buông thõng xuống sàn: ``Sao tự nhiên hôm nay trời nóng dữ vậy trời\dots''

Flare ngửa mặt lên trần nhà, mồ hôi rịn ra trên trán: ``Mới chỉ vào xuân thôi mà thời tiết đã thế này rồi\dots''

Towa thè lưỡi: ``Ước gì bây giờ có cái gì đó thật mát để giải khát nhỉ\dots''

Flare gật đầu đồng tình: ``Đồng ý hai tay hai chân luôn\dots''

*CẠCH*

Tiếng cửa phòng bật mở. Kuon bước vào, theo sau Kaigou cùng Naomi. Cậu Tổng vừa cởi chiếc mũ lưỡi trai xuống vừa quạt quạt lấy hơi, trong khi nàng Tổng một tay xách chiếc thùng giữ nhiệt lớn chứa đầy đá lạnh một cách nhẹ tênh như cầm túi bông gòn. Bên cạnh cô, bé Naomi ngoan ngoãn ôm một xấp cốc nhựa và ống hút, đôi mắt xanh lục bảo trong veo ngơ ngác nhìn quanh.

Kuon lên tiếng: ``Chào buổi sáng mấy đứa. Hôm nay trời nóng kinh khủng nhỉ?''

Hai cô gái trên sofa uể oải đến mức còn không buồn ngồi dậy, chỉ biết ú ớ vài tiếng gợn sóng: ``Nóng quá ạ\dots'' và ``Khát muốn chết mất thôi\dots''.

Kaigou đặt chiếc thùng giữ nhiệt cái *CỘP* xuống mặt bàn, khoanh tay bĩu môi nhìn hai cô nàng đang nằm dài: ``Mới vào xuân oi một tí mà hai đứa đã nhũn ra như bánh đa ngâm nước thế kia rồi à? Onii, mau chia nước cho họ đi kẻo cái văn phòng này biến thành phòng cấp cứu bây giờ.''

Naomi lon ton bước lại gần bàn tròn, cẩn thận đặt xấp cốc nhựa xuống rồi ngước đôi mắt tròn xoe lên nhìn Towa và Flare, cất giọng trong trẻo nhưng đầy vẻ điềm đạm: ``Towa-oneechan, Flare-oneechan, hai chị ráng một chút nhé. Em với Onii-chan và Kaigou-oneechan có mang thức uống mát lành tới nè\~{}''

Nói rồi, Kuon mở nắp thùng giữ nhiệt. Kaigou thoăn thoắt nhấc ra ba chiếc cốc lớn chứa đầy thứ nước màu vàng óng ả ngả xanh lá tươi mát, bên trên nổi lớp bọt mịn màng cùng những lát quất thái mỏng và đá viên lạnh buốt tỏa ra làn sương mờ ảo.

Chẳng hiểu lấy đâu ra sức mạnh, ngay khi ba ly nước vừa đáp xuống bàn, cả Towa và Flare liền bật dội dậy như có lò xo dưới mông, nhanh chóng sán lại gần quan sát.

Towa nheo mắt nhìn ly nước màu sắc lạ lẫm: ``Trông món này lạ thế nhỉ\dots''

Flare gật gù, ngửi thấy một mùi thơm thanh mát dịu ngọt: ``Nhìn cũng có vẻ rất giải khát nữa!''

Towa thắc mắc quay sang hỏi: ``Cơ mà\dots\ ly nước này khác gì so với trà sữa hay nước ngọt bọn em hay uống vậy anh Kuon?''

Kuon thong thả giải thích: ``Khác nhiều chứ em. Đây là 'Nước mía', thức uống giải khát quốc dân cực kỳ phổ biến ở quê anh. Người ta ép trực tiếp từ những cây mía tươi nguyên chất rồi vắt thêm chút quất để dậy mùi thơm thanh dịu.''

Kaigou tựa lưng vào bàn, khoanh tay tiếp lời với vẻ đắc ý: ``Đường trong ly này hoàn toàn là đường tự nhiên từ thân cây mía, vừa nạp năng lượng tức thì lại vừa thanh nhiệt, tốt cho sức khỏe gấp vạn lần mấy cái siro công nghiệp trong trà sữa của mấy đứa đấy. Ở quê chị trưa hè mà làm một ly to đùng thế này thì có chạy bộ chục cây số cũng chẳng biết mệt là gì.''

Naomi đưa từng chiếc ống hút cho Towa và Flare, khẽ nghiêng đầu nói nhỏ: ``Vị ngọt của đất trời Etsunan đó ạ\dots\ Naomi nếm thử một ngụm lúc ép mía mà thấy tâm hồn thanh thản ạ.''

Flare gật gù tâm đắc: ``Nếu nó vừa mát lại vừa tốt cho sức khỏe hơn trà sữa thì\dots''

Towa lanh chanh chộp lấy ống hút: ``\dots Bọn em phải uống ngay mới được!''

Cả hai vội vàng cắm ống hút vào ly, hít một hơi thật sâu và thưởng thức ngụm đầu tiên.

*SỤT\dots!*

Ngay khoảnh khắc dòng nước mía mát lạnh tuôn chảy qua cổ họng, một cảm giác sảng khoái đến tận cùng bùng nổ. Trong tâm trí của Towa và Flare, họ như đang thấy mình đứng giữa một bạt ngàn đồng mía xanh tươi rợp bóng mát, nơi những cơn gió cao nguyên thổi qua đánh tan hoàn toàn cái nóng oi ả ban nãy.

``\textbf{Mát quáaaa!!}'' cả hai đồng thanh reo lên đầy sung sướng.

Không chút chần chừ, hai cô gái bắt đầu hút sùn sụp liên tục. Chỉ trong chưa đầy nửa phút, hai ly nước mía lớn đã sạch bách không còn một giọt, chỉ còn lại những viên đá nhỏ dưới đáy.

Kuon đứng bên cạnh trố mắt nhìn tốc độ tiêu thụ kinh hoàng của hai người: ``H-Hai đứa uống nhanh thế!? Bộ cảm thấy nóng đến mức đó sao?''

Towa và Flare đồng thanh đáp, mắt long lanh nhìn vào thùng giữ nhiệt: ``Trời hôm nay nóng thật mà anh!! Còn ly nào nữa không, hai anh chị!?''

Kuon đổ mồ hôi hột: ``Nào nào, cứ bình tĩnh. Thùng này Kaigou mang theo nhiều lắm, cứ uống thoải mái đi.''

``\textbf{Tuyệt vời quá!!}''

Kaigou phì cười, vỗ mạnh lên vai Kuon: ``Thấy chưa Aniki? Em đã bảo là mang một thùng lớn đi không bao giờ thừa mà lại. Nào, để chị lấy thêm cho.''

Và rồi, điều gì đến cũng phải đến. Đúng lúc đó Aki và Ayame từ ngoài đi vào, thấy món nước giải nhiệt thần thánh liền gia nhập cuộc vui. Bằng một phép thuật nào đó, chỉ với bốn cô gái \hll\ mà cả cái thùng giữ nhiệt to đùng của Kaigou mang theo đã bị `vét sạch' không còn một giọt!

Cậu Tổng nhìn đống ly rỗng trên bàn mà mặt xám giét, lo sốt vó: ``Kiểu này chắc sắp tới anh phải mua hẳn một cái máy ép nước mía cùng hàng chục cây mía tươi đặt ở văn phòng mới đủ phục vụ mấy đứa quá\dots''

Kaigou cười khẩy, huých nhẹ tay anh trai: ``Anh cứ mua về đi Aniki, em tình nguyện đứng quay tay ép mía cho, bảo đảm ép kiệt đến giọt cuối cùng không sót một tí nước nào!''

Trong khi mọi người đang cười đùa rôm rả, cô bé Naomi bỗng nghiêng đầu, đôi mắt xanh lục bảo chăm chú nhìn bốn cô gái đang ngồi xoa bụng thỏa mãn.

Cô bé chớp chớp mắt, cất giọng hỏi ngây thơ: ``À\dots\ Naomi hỏi chút\dots\ Chiều tối nay Aki-oneechan, Ayame-oneechan với Towa-oneechan có buổi phát sóng trực tiếp hay biểu diễn gì không ạ?''

Towa gật đầu tắp lừ: ``Có chứ! Tối nay chị với Ayame-senpai có lịch stream collab game chung nè.''

Aki cũng mỉm cười: ``Chị cũng có buổi stream nói chuyện với fan nữa.''

Nghe đến đó, nụ cười trên môi Kaigou bỗng khựng lại. Cô chớp mắt hai lần, sắc mặt hơi biến đổi rồi quay ngoắt sang nhìn Kuon: ``Khoan đã\dots\ Aniki. Anh có nhớ ở giới nghệ sĩ biểu diễn và sân khấu quê mình có một cái điều kiêng kỵ dân gian gì không?''

Kuon giật thót mình, mắt mở to: ``Chết dở\dots\ đừng nói là vụ kiêng cữ đó nhé!?''

Cả bốn cô gái tròn mắt tò mò: ``Là gì vậy, Kaigou-senpai?''

Kaigou thì thầm đầy bí hiểm: ``Giới nghệ thuật truyền tai nhau một quy tắc tâm linh tối kỵ: tuyệt đối không được uống nước mía trước giờ lên sóng hay biểu diễn! Người ta tin rằng nước mía ngọt ngào nhưng thân mía chia từng đốt rỗng ruột, uống vào là y như rằng dễ bị 'bể show', mất mạng, sập kênh, hoặc gặp trục trặc kỹ thuật kinh hoàng giữa chừng đấy!''

Bé Naomi gật đầu phụ họa, nét mặt nghiêm túc hệt như một bà cụ non: ``Dạ đúng rồi ạ\dots\ Các cụ bảo là vía cây mía dễ làm mọi chuyện vỡ vụn ra từng khúc, không trọn vẹn được đâu ạ.''

Vừa nghe xong câu đó, cả Towa, Flare, Aki và Ayame đều hóa đá tại chỗ, da gà da vịt nổi rần rần khắp cánh tay.

Kuon đứng bên cạnh cũng nín thở, mồ hôi hột chảy dài trên má nhìn bốn cô gái đang run rẩy vì lo cho buổi stream tối nay: ``Hy vọng\dots\ hy vọng là tâm linh Etsunan không `ứng nghiệm' sang tận trụ sở COVER Corp ở Nhật Bản này!''

\newpage

\subsection*{\phantomsection\label{ep-C9.V}}
\addcontentsline{toc}{subsection}{{\sen{-Histoire 5-}} {\caviar Sa tế tôm}}
\stpt{-Histoire 5-}
\stcap{Sa tế tôm}
\vspace{-4pt}
\begin{center}
  {\caviar\fontsize{12pt}{12pt}\selectfont Nhân vật:}
\end{center}

Một lần nọ, Aki tức tốc xông vào văn phòng \hll, hào hứng gọi lớn: ``Mọi người ơi! Mình vừa xem được cái này hay lắm, phải hỏi mọi người ngay mới được!''

Lời kêu gọi đầy năng lượng của cô lập tức thu hút sự chú ý của Ayame và Towa. Ở góc phòng, Kuon vẫn đang gõ máy tính dở tay, còn Ryuuhou thì đang ngồi bên cạnh phụ anh trai sắp xếp lại chồng tài liệu kịch bản.

Aki tiến lại gần bàn tròn, mắt lấp lánh: ``Hôm vừa rồi trên TV, chị xem được một loại gia vị này nhìn độc lạ lắm!''

Nàng Ác ma tò mò rướn người tới: ``Gia vị gì thế, Aki-senpai?''

Nàng Dị giới xòe hai tay miêu tả: ``Chả là\dots\ Nó là một loại sốt dạng dầu cay, thấy người ta bảo dùng để ăn kèm với cơm trứng sống (Tamago Kake Gohan) ngon cực kỳ!''

Nàng Oni nghiêng đầu: ``Trông nó thế nào hả chị?''

``Xem nào\dots\ Nó có màu đỏ sẫm, hơi sền sệt, bên trong có mấy cái vảy nhỏ nhỏ\dots''

Kuon dừng tay gõ bàn phím, im lặng lắng nghe.

``Mùi hương thơm lừng, được miêu tả là có vị cay nồng bùng nổ\dots\ lại còn được đựng trong một cái lọ thủy tinh nữa!''

Towa mất kiên nhẫn: ``Rốt cuộc tên nó là gì vậy, Aki-senpai?''

Nàng Dị giới đưa ngón tay lên cằm, đắn đo hồi lâu rồi ấp úng phát âm: ``Chị\dots\ chị không nhớ rõ lắm, nhưng hình như nhãn hiệu ghi là\dots\ `sa-tê-tôm-mư'?''

Hai cô nàng còn lại nghe cái tên lạ hoắc thì ngơ ngác nhìn nhau.

``Sa-tê-tôm-mư?'' Towa nhíu mày.

``Tên nghe kỳ lạ thật đấy-yo.'' Ayame gật gù.

Aki gãi đầu phân trần: ``Chị cũng không biết nữa, nhãn chai ghi chữ Latinh nhưng đọc vào chẳng giống tiếng Anh chút nào cả!''

Cậu Tổng phì cười, cất tiếng từ phía bàn làm việc: ``Đó là Sa tế tôm của Etsunan chứ gì nữa.''

Cả nhóm giật mình quay lại. ``Sa tế\dots\ Ể? Cái gì cơ?''

Ryuuhou mỉm cười giải thích hộ anh trai: ``Đó là một loại sốt ớt chưng đặc trưng ở quê em đấy ạ! Người ta làm nó từ ớt bột, sả bào mỏng, tôm khô giã nhỏ rồi chưng ngập trong dầu thơm phức. Nhưng mà\dots''

Cô quay sang nhìn anh trai, hai anh em cùng dở khóc dở cười. Kuon lắc đầu: ``Anh cũng thấy lạ đấy. Bên quê anh, sa tế tôm chủ yếu dùng để nêm nước lẩu hoặc ướp đồ nướng. Chưa thấy ai lại đi chan sa tế tôm vào cơm trứng sống bao giờ cả!''

Aki chớp mắt phân sưa: ``Ể? Nhưng trên chương trình ẩm thực TV người ta ăn thử rồi khen nức nở lắm mà?''

Kuon nhún vai: ``Thì mỗi người một khẩu vị. Thôi được rồi, nếu mấy đứa tò mò thì cứ thử xem sao.''

Chẳng mấy chớp mắt, nàng Dị giới đã chạy tót ra cửa hàng tiện lợi mua về một lọ sa tế tôm. Cô xới sẵn bốn bát cơm nóng hổi, mỗi bát đập một quả trứng gà sống béo ngậy lên trên.

``Bây giờ mỗi người tự cho lượng sa-tê-tôm-mư vừa ăn vào bát của mình nha!'' cô hào hứng chuyền tay lọ sa tế.

Aki và Towa cẩn thận múc khoảng một thìa cà phê vừa đủ. Ayame thì bạo dạn múc lấy múc để, ``bay hơi'' mất gần một phần ba lọ sa tế màu đỏ quạch.

Đến lượt Kuon, cậu chỉ khẽ dùng đầu đũa gắp một tẹo vảy sa tế và rưới vài giọt dầu ớt lên trên. Ryuuhou nhìn bát của anh trai mình rồi tủm tỉm cười. Mọi người trong phòng ai nấy đều thừa biết cấp độ ăn cay của cậu Tổng ``bá đạo'' đến mức nào.Với một người có thể nhai ớt chỉ thiên ngọt xáng như ăn ớt ngọt như Kuon, việc cậu chỉ gắp vài hạt sa tế tượng trưng thế này hoàn toàn không phải vì sợ cay, mà đơn giản là vì cậu thấy sự kết hợp sa tế với trứng sống nó\dots\ ``sai sai'' quá!

``Mời mọi người ăn cơm! / Mời cả nhà ăn cơm-yo!!''

Cả bốn người cùng trộn đều và xúc một thìa lớn cho vào miệng. Phản ứng lập tức chia thành những thái cực rõ rệt.

Aki và Towa vừa nuốt xong liền chắp tay, mắt sáng quắc như nhìn thấy ánh hào quang từ thiên đường chiếu xuống. Sự kết hợp giữa vị béo ngậy của trứng sống với độ đậm đà, thơm lừng mùi sả tôm và cay nhẹ của sa tế tạo nên một hương vị bùng nổ ngoài mong đợi.

Nàng Dị giới ngất ngây: ``Vị cay the the thơm nức của sa-tê-tôm-mư\dots''

Nàng Ác ma tiếp lời: ``\dots kết hợp với sự mềm béo ngậy của trứng sống\dots''

Cả hai đồng thanh reo lên: ``\textbf{Thật là tuyệt hảo quá đi mất!!}''

Kuon đổ mồ hôi hột nhìn hai người: ``Hai đứa mê mẩn món này thật đấy à? Còn Ayame-chan thì sa--''

Chưa kịp để cậu nói hết câu, nàng Quỷ sừng đã há miệng\dots\ hà ra một luồng khí nóng như sắp phun ra lửa! Dù mặt đỏ gay và hai tai bốc khói vì múc quá tay, cô nàng vẫn phấn khích đập bàn: ``Hộc\dots\ Hộc\dots\ Cay thật đấy-yo! Nhưng mà ngon dã man! Cái này mà đem lên livestream thách thức ăn cay chắc chắn fan sẽ thích lắm đây!''

Ryuuhou nếm thử một thìa nhỏ trong bát của mình rồi gật gù, quay sang nhìn Kuon: ``Cũng không tệ lắm anh nhở? Nhưng đúng là ăn thế này nó cứ thiếu thiếu cái gì ấy\dots''

Kuon thong thả nhai nốt miếng cơm, điềm nhiên nhận xét: ``Ừ, không tệ, nhưng vị nó không hợp rơ hoàn toàn. Sa tế tôm sinh ra là để dành cho những món nóng hổi bốc khói cơ.''

Aki tò mò: ``Là món gì vậy?''

Kuon đáp: ``Ở quê anh, sa tế tôm phải thả vào mấy nồi lẩu nghi ngút khói, hoặc ít nhất là dùng để xào cơm rang sa tế. Vị nóng của nước dùng nó mới đẩy được hết cái thơm của sả và tôm khô lên.''

Towa chép miệng: ``Ồ\dots\ Ra là vậy!''

Ayame sáng mắt lên: ``Lẩu ăn với sa tế nghe hấp dẫn quá-yo!''

Aki lập tức đề xuất: ``Nếu đúng như anh nói, sao chúng mình không thử làm một nồi lẩu ngay bây giờ luôn nhỉ?''

Kuon kinh ngạc nhìn ba cô gái: ``Mấy đứa muốn ăn lẩu ngay lúc này ư?'' ba người gật đầu lia lịa.

``Giữa cái thời tiết nắng bốc hỏa này sao!?'' cả ba lại gật đầu tắp lừ không chút do dự.

Kuon bật cười đầu hàng trước tâm hồn ăn uống của họ: ``Được rồi, thích thì chiều! Thế mấy đứa muốn ăn lẩu gì?''

Cuối cùng, cả nhóm quyết định chốt đơn một nồi lẩu Shabu-shabu hoành tráng, và tất nhiên không thể thiếu lọ sa tế tôm làm linh hồn cho bát nước chấm.

Phải công nhận một điều, khi nhúng miếng thịt bò thái mỏng vào nước lẩu sôi ùng ục rồi chấm ngập trong bát sốt sa tế tôm béo ngậy, cay nồng, cả năm người - kể cả Kuon - đều phải gật gù thừa nhận: Món này đúng quả là ``số dzách''!

\newpage

\subsection*{\phantomsection\label{ep-C9.VI}}
\addcontentsline{toc}{subsection}{{\sen{-Histoire 6-}} {\caviar Chè hạt sen}}
\stpt{-Histoire 6-}
\stcap{Chè hạt sen}
\vspace{-4pt}
\begin{center}
  {\caviar\fontsize{12pt}{12pt}\selectfont Nhân vật: \emp{naomi2} \emp{aki} \emp{ayame} \emp{flare1} \emp{towa}}
\end{center}

Một ngày khác tại văn phòng \hll.

Trên chiếc bàn tròn ở giữa phòng nghỉ đột nhiên xuất hiện một chiếc lon nhôm bí ẩn. Lon hoàn toàn trọc lóc, không dán nhãn mác, không có bất kỳ dòng chữ hướng dẫn nào, chỉ mang thiết kế nắp giật đơn thuần giống như mấy lon nước ngọt thông thường.

Hôm nay Kuon và Ryuuhou có lịch trình riêng nên không đến văn phòng, còn Kaigou thì đang bận việc bên nhánh \hls. Nơi này chỉ có Ayame, Aki, Towa và bé Naomi đang ngoan ngoãn ngồi xếp giấy origami ở góc bàn tròn.

Ba cô gái trẻ đứng vây quanh chiếc lon kim loại, ánh mắt tràn đầy sự nghi ngại.

Aki chỉ tay: ``Trên bàn kia là chiếc lon gì thế hai người?''

Ayame lắc đầu: ``Em chịu thôi. Lúc em đến nó đã nằm chễm chệ ở đấy rồi. Chắc lại là Towa-sama bày trò gì hãm hại chị em mình đúng không?''

Towa giật nảy mình, xua tay lia lịa: ``K-Không phải em! Em còn tưởng Aki-senpai hay Ayame-senpai mang tới chứ!''

Aki quay sang nhìn cô bé út nhà Hikawa, ngập ngừng hỏi: ``Naomi-chan\dots\ chiếc lon này không ghi tên tuổi gì cả. Em có biết ai mang tới đây không?''

Naomi dừng tay gấp hạc giấy, ngước đôi mắt xanh lục bảo trong veo lên nhìn chiếc lon nhôm, rồi khẽ lắc đầu cất giọng trong trẻo: ``Naomi không rõ của ai mang tới ạ. Lúc em vào phòng đã thấy nó nằm đó rồi. Nhưng mà\dots\ Naomi ngửi thấy mùi thơm rất thanh khiết và ngọt dịu của hoa sen, không có khí tức độc hại hay nguy hiểm gì đâu, Aki-oneechan.''

Dù được cô bé thấu thị trấn an, ba cô gái trẻ vẫn còn chút dè chừng. Thế là cả ba đứng khoanh tay, chụm đầu quan sát cái lon suốt nửa tiếng đồng hồ. Không ai dám động vào.

Cuối cùng, không chịu nổi sự tò mò thiêu đốt, nàng Quỷ sừng Ayame dũng cảm bước lên: ``Thôi được rồi! Để em ra mở thử xem sao-yo!''

Towa lo lắng níu tay: ``Kìa Ayame! Nhỡ bên trong là một quả bom nổ chậm thì sao!?''

Aki cũng can ngăn: ``Nhỡ mở ra có con gì đáng sợ bò ra thì tiêu đấy!''

Ayame vỗ ngực trấn an: ``Chắc không đến mức tệ thế đâu mà\dots''

*PÓC!*

Tiếng nắp nhôm bật mở vang lên nhẹ nhàng. Không có khói bốc ra, cũng chẳng có quả bom nào nổ. Ayame ghé mắt nhìn vào bên trong: ``Ê-tồ\dots\ Cái gì thế này?''

Bên trong lon là một thứ nước sóng sánh, thơm nhẹ. Phía trên nổi lên mấy hạt tròn màu vàng ươm, cùng vài quả tròn màu trắng trong suốt chìm bên dưới.

Aki hít hà không khí xung quanh chiếc lon: ``Hình như\dots\ mình ngửi thấy đúng mùi hạt sen như Naomi-chan vừa nói kìa!''

Towa gật gù: ``Em cũng ngửi thấy thế. Mùi hương dễ chịu quá\dots''

Ayame nhanh tay trút toàn bộ thức uống trong lon ra một chiếc bát thủy tinh to gần đó. Cả ba cô gái lại gần quan sát kỹ hơn.

``Nhìn giống một món chè ngọt\dots'' Aki nhận định. ``Mấy hạt màu vàng này là hạt sen đúng không nhỉ?''

``Chắc chắn rồi,'' Towa đáp. ``Còn mấy quả màu trắng ngà này chắc là\dots''

``Nhãn đúng không?'' Ayame tiếp lời, rồi quay sang hỏi hai người kia: ``Có ai muốn ăn thử không?''

Aki và Towa vẫn còn chút dè chừng. Nhưng thấy món ăn tỏa ra mùi thơm dịu nhẹ rất tự nhiên, thanh khiết nên\dots

``Chắc\dots\ chắc là ăn được, không sao đâu nhỉ?''

Ayame xung phong múc một thìa chè đầy hạt sen đưa vào miệng đầu tiên. Ngay lập tức, đôi mắt cô nàng sáng rực lên, thốt lên đầy kinh ngạc: ``Ngọt thanh ngon quá đi mất-yo!''

Thấy Ayame hoàn toàn bình an vô sự lại còn khen nức nở, Aki liền khéo léo múc một quả nhãn mọng nước, còn Towa thì múc một thìa đầy hạt sen lẫn nước chè ngọt mát. Cả hai cùng nếm thử.

Ngay khoảnh khắc vị ngọt thanh dịu tan trên đầu lưỡi, cả Aki và Towa như rơi vào không gian thư thái giữa một hồ sen thơm ngát.

Towa sướng rên: ``Hạt sen ninh bở bùi bùi, ăn sần sật thích thật đấy\dots!''

Aki gật gù ngất ngây: ``Cộng với vị nhãn mọng nước ngọt dịu nhẹ kết hợp với nước chè thanh mát\dots!''

Cả hai đồng thanh reo hò: ``Phải nói là ngon không còn gì để chê luôn!''

Ayame nhìn hai người bạn khen ngợi có phần hơi ``thái quá'' mà đổ mồ hôi hột, nhưng rồi cũng mỉm cười mãn nguyện. Cứ thế, ba cô gái hào hứng chia nhau ăn sạch bách bát chè hạt sen long nhãn, không còn sót lại dù chỉ một giọt nước chè hay một hạt sen!

Bé Naomi ngồi một góc nhìn ba chị lớn tranh nhau ăn ngon lành, khóe môi khẽ cong lên một nụ cười nhẹ nhõm, đôi mắt xanh lấp lánh vẻ hồn nhiên.

Đúng lúc cả ba đang xoa bụng mãn nguyện thì cửa văn phòng bật mở. Flare bước vào, vừa đi vừa ngân nga giai điệu vui tươi: ``Trưa nắng thế này mà được thưởng thức lon chè hạt sen long nhãn ướp lạnh mình mới mua thì còn gì bằng\~{}``

Thế nhưng, vừa nhìn thấy chiếc lon nhôm tróc nhãn nằm trần truồng trên bàn, bên cạnh là cái bát thủy tinh rỗng tuếch, Flare đứng hình mất năm giây. Cô nàng Yêu tinh tóc vàng hét lên trong sự sốc nặng: ``\textbf{Áaaaaa!! Ai mở lon chè hạt sen của tôi vậy!?? Ai ăn hết rồi!??}''

Trong phòng nghỉ chỉ còn lại Flare đang vò đầu bứt tóc trong cơn uất ức, và bé Naomi vẫn điềm nhiên ngồi gấp nốt cánh hạc giấy cuối cùng.

Naomi đặt chú hạc giấy xuống bàn, nhẹ nhàng bước lại gần, kéo nhẹ vạt áo Flare. Cô bé ngước đôi mắt to tròn trong veo lên nhìn đàn chị, cất giọng non nớt nhưng an ủi đầy chân thành: ``Flare-oneechan đừng buồn nhanha\dots\ Mấy chị ban nãy ăn xong khen ngon lắm đó ạ. Để tối về em nhờ Suzuno-oneechan với Onii-chan nấu một nồi chè hạt sen long nhãn thật to cho chị ăn bù nhé?''

Nhìn thấy gương mặt ngây thơ, ngoan ngoãn và ánh mắt thấu hiểu của cô bé út nhà Hikawa, ngọn lửa giận trong lòng Flare bỗng chốc tan biến sạch ghẽ.

Cô nàng thở dài thượt, dở khóc dở cười cúi xuống xoa xoa mái tóc ngắn của Naomi:  ``Trời ơiơi\dots\ em gái út nhà Hikawa ngoan thế này thì làm sao chị nỡ giận cho được chứ\dots\ Coi như chị đãi mấy đứa quậy phá kia vậy!''

\newpage

\subsection*{\phantomsection\label{ep-C9.VII}}
\addcontentsline{toc}{subsection}{{\sen{-Histoire 7-}} {\caviar Mắm tôm}}
\stpt{-Histoire 7}
\stcap{Mắm tôm}
\vspace{-4pt}
\begin{center}
  {\caviar\fontsize{12pt}{12pt}\selectfont Nhân vật: \emp{kuon2} \emp{ryuuhou2} \emp{aki} \emp{ayame} \emp{flare1} \emp{towa}}
\end{center}

Một không khí hỗn loạn chưa từng có đang diễn ra tại phòng nghỉ văn phòng \hll. Cả bốn cô gái - Aki, Flare, Ayame và Towa - đang hoảng loạn chạy nháo nhào khắp phòng, hai tay bịt chặt mũi vì một mùi hương cực kỳ ``kinh dị'' bất ngờ giăng lối khắp không gian.

Towa gào lên, hai tai quỷ quăn lại: ``\textbf{C-Cái mùi thối gì thế này!?}''

Aki vừa chạy vừa sặc sụa: ``A\textbf{i đó mở cửa sổ ra mau! Không khí độc mất thôi!}''

Ayame nhăn nhó, mắt rơm rớm nước mắt: ``Hộc\dots\ hộc\dots\ Em không thể thở nổi với cái mùi này đâu, yo!''

Trong khi ba người kia đang mất phương hướng, cô nàng Yêu tinh tóc vàng Flare với chiếc mũi thính nhạy lại vô tình phát hiện ra nguồn gốc. Cô run rẩy chỉ tay về phía cánh cửa tủ bếp đang mở hé ở góc phòng: ``M-Mọi người! Cái mùi này\dots\ nó phát ra từ trong tủ bếp kia kìa!''

Nghe vậy, ba người còn lại lập tức tiến lại gần hướng chỉ tay của Flare. Nhưng càng tiến lại gần chiếc tủ, mùi hôi nồng nặc ma quái càng xộc thẳng vào mũi dữ dội hơn.

``K-Không ổn rồi! Mùi nặng quá!'' nàng Quỷ sừng thảng thốt lùi lại.

Và rồi, sự cố tồi tệ đầu tiên đã xảy ra. Nàng Ác ma Towa do hít phải một ngụm khí quá đậm đặc đã trợn ngược mắt, đổ gục xuống sàn nhà, hai mắt quay mòng mòng như chong chóng.

``T-Towa-sama!? Tỉnh lại đi em!'' Aki hốt hoảng buông tay khỏi mũi để lay người hậu bối.

Thế nhưng, chính vì buông tay ra đúng lúc ngọn gió mang theo ``luồng khí độc'' tràn qua, cô ngay lập tức hóa đá, mắt nhắm tịt rồi ``đo sàn'' ngay bên cạnh Towa.

Flare thót tim kêu lên: ``\textbf{Aki-senpai ngất luôn rồi kìa! Trời ơi!}''

Ayame nắm chặt tay cô, ánh mắt kiên định như những chiến binh ra trận: ``Không còn cách nào khác đâu Flare! Chúng ta phải dũng cảm tiến lên đóng cánh cửa tủ đó lại để cứu lấy văn phòng này!''

Gạt nước mắt bỏ lại hai người đồng đội đã bị ``hạ đo ván'', hai người nín thở, thận trọng từng bước nhích lại gần chiếc tủ bếp. Sau một hồi giằng co căng thẳng, cả hai cuối cùng cũng chạm được tay vào mép cửa tủ.

``Đếm đến ba rồi cùng dùng lực đóng sầm nó lại nhé!'' Flare thì thầm.

``Được-yo!''

``Một\dots\ hai\dots\ \textbf{ba!!}''

Thế nhưng, trong sự cuống quýt và hoảng loạn tột độ, thay vì dùng lực đẩy cánh cửa tủ đóng lại, cả hai lại vô tình\dots\ kéo toang cánh cửa tủ ra hết cỡ.

Một hũ mắm tôm màu tím thẫm khổng lồ đang mở nắp trần trụi nằm ngay trung tâm chiếc tủ. Luồng ``khí độc'' nén bấy lâu nay bùng nổ, sộc thẳng trực diện vào mặt hai cô gái.

``Chết toi rồi! Nhầm hướng rồi--!''

Đó là những lời trăn trối cuối cùng của Flare và Ayame trước khi cả hai mắt nhắm tịt, đổ gục xuống đất nằm la liệt bên cạnh Aki và Towa.

Vài phút sau, cánh cửa phòng nghỉ mở ra. Kuon bước vào với một đĩa bún đậu mắm tôm nóng hổi, giòn rụm trên tay. Đi ngay bên cạnh cậu là cô em gái Ryuuhou, tay xách theo túi rau kinh giới và dưa chuột tươi rói.

Vừa bước vào phòng, đập vào mắt hai anh em là cảnh tượng bốn cô gái \hll\ đang nằm chần ngần, la liệt trên sàn nhà như một ``chiến trường'' thực sự.

Cậu Tổng giật mình trố mắt: ``Mấy đứa làm sao thế này!? Có chuyện gì xảy ra vậy!?''

Nàng Trưởng nhóm nhìn quanh một lượt, rồi ánh mắt cô bé dừng lại ở cánh cửa tủ bếp đang mở toang. Cô em gái 15 tuổi lập tức hiểu ra vấn đề, khẽ vỗ trán phì cười: ``Trời ạ! Em đã dặn Onii là lấy hũ mắm tôm ra pha chế xong thì phải vặn chặt nắp cất đi cơ mà! Anh lại quên đóng nắp tủ đúng không?''

Kuon gãi đầu cười trừ đầy tội lỗi: ``À\dots\ Hè hé, anh tiện tay để đó tí nữa cất, ai ngờ\dots''

Lúc này, Flare từ từ chống tay ngồi dậy, đầu óc vẫn còn choáng váng: ``M-Mắm tôm sao\dots\ Kuon-senpai\dots''

Ayame lè lưỡi thào thào: ``Hôi quá\dots\ Muốn ngất lần hai luôn-yo\dots''

Aki và Towa chắp tay mếu dào: ``Chịu hết nổi rồi\dots\ Cứu bọn em với\dots''

Ryuuhou bước lại gần, vừa nhanh tay vặn chặt nắp hũ mắm tôm cất đi vừa trêu chọc: ``Mấy chị yếu bóng vía thế! Mắm tôm Etsunan nhà em tuy mùi hơi nồng tí thôi, nhưng pha thêm chút đường, chanh, ớt rồi đánh bông lên chấm bún đậu thì là `cực phẩm' luôn đấy! Onii làm cho bọn em ăn suốt mà!''

Thế nhưng, bất chấp lời giải thích đầy hào hứng của Ryuuhou, báo hại là suốt cả ngày hôm đó, bốn cô gái \hll\ đã kiên quyết tránh xa Cậu Tổng cả chục mét, chỉ vì cả người cậu đã bị ám trọn cái mùi ``mắm tôm huyền thoại'' ấy!

\newpage

\subsection*{\phantomsection\label{ep-C9.VIII}}
\addcontentsline{toc}{subsection}{{\sen{-Histoire 8-}} {\caviar Trứng vịt lộn}}
\stpt{-Histoire 8-}
\stcap{Trứng vịt lộn}
\vspace{-4pt}
\begin{center}
  {\caviar\fontsize{12pt}{12pt}\selectfont Nhân vật: \emp{kuon2} \emp{ryuuhou2} \emp{suzuno2} \emp{aki} \emp{ayame} \emp{flare1} \emp{towa}}
\end{center}

Một buổi chiều nọ, Ayame hứng khởi tổ chức một cuộc thi ăn trứng ngay tại bàn tròn văn phòng.

``Chào mừng mọi người đến với Cuộc thi Ăn Trứng! Luật chơi rất đơn giản: ai ăn được nhiều trứng nhất trong thời gian quy định sẽ giành chiến thắng!'' Ojou dốc sức hô hào.

Thế nhưng, sau lưng mọi người, nàng Quỷ sừng lại thầm thì cười khoái chí: ``\textit{Khì khì khì\dots\ Không ai biết được đây thực chất là một cái bẫy đâu! Trong đĩa trứng này, ngoài trứng luộc bình thường, mình đã lén trộn vào vài quả `trứng vịt lộn' siêu kinh dị mà mình mua được ở siêu thị đồ ngoại nhập! Lần này mình sẽ dọa họ một trận khiếp vía cho xem!}''

Tất nhiên, với sự tinh ý bẩm sinh, một số người đã lập tức đánh hơi thấy sự bất thường.

Flare nheo mắt nhìn nụ cười gian xảo của Ayame: ``Không biết Ayame-senpai lại đang bày trò ma mãnh gì đây\dots''

Towa khoanh tay nghi ngờ: ``Xem chừng có vẻ mờ ám lắm nha\dots''

Thế nhưng, cũng có những tâm hồn ăn uống hoàn toàn không bận tâm đến cạm bẫy.

Aki ôm bụng thều thào: ``Món gì cũng được\dots\ Chị đói quá, từ sáng tới giờ chưa có gì bỏ bụng cả\dots''

Suzuno dịu dàng ngồi bên cạnh Aki, gật đầu đồng tình: ``Đúng đó, em cũng đang thèm chút gì đó béo béo ngậy ngậy nè.''

Trong khi đó, Kuon và Ryuuhou lại tỏ ra vô cùng thong dong. Cô em gái chống cằm cười khẩy nhìn Ayame, còn ông anh trai thì gật gù: ``Xem chừng cuộc thi này có vẻ thú vị đấy.''

Thấy các thí sinh đã tập hợp đầy đủ, Ayame vỗ tay ròn rã: ``Các thí sinh sẵn sàng\dots\ Chuẩn bị\dots\ \textbf{Bắt đầu!}''

Ngay lập tức, Aki, Flare, Towa, Kuon, Suzuno và Ryuuhou đều với tay lấy quả trứng đầu tiên. Kuon, Aki và Flare bóc vỏ ra một quả trứng gà luộc chín vàng ươm bình thường.

Thế nhưng, ``vận may'' lại gọi tên\dots\ nàng Ác ma Towa!

*CẠCH!*

Ngay khi Towa gõ nhẹ quả trứng xuống đĩa và bóc ra một mảng vỏ, đôi mắt cô nàng trố ngược. Cô thốt lên một tiếng hét thất thanh làm cả văn phòng nhảy dựng lên: ``\textbf{Á-áaaaa!! Có con gì ở bên trong này!!}''

Trên đĩa của cô, phần vỏ vỡ ra để lộ một cấu trúc kỳ dị: một nửa vàng cam với những đường mạch máu chăng chịt, và nửa còn lại là một hình hài sinh vật nhỏ bé đầy lông tơ nhạt!

Towa xua tay nhảy tót lên ghế: ``C-c-con gì thế này!?? Đ-Đáng sợ quá!!''

Aki giật mình che miệng: ``Trông kinh khủng thế!''

Flare cũng toát mồ hôi hột: ``C-Cái này là sinh vật gì vậy!?''

Trái ngược hoàn toàn với sự hoảng loạn của ba cô gái \hll, ba anh em nhà Hikawa lại có phản ứng bình thản đến mức kỳ quặc.

Kuon điềm nhiên dùng thìa xúc một ngụm nước ngọt từ trong quả trứng của mình, đáp nhẹ hẫng: ``À, trứng vịt lộn ấy mà.''

Đến cả Ayame - ``thủ mưu'' của trò đùa này - cũng phải ngơ ngác trước sự bình tĩnh của Cậu Tổng: ``A-Anh biết món này sao, Kuon-senpai?''

Cậu gật đầu: ``Ừ, đặc sản quê anh mà. Trứng vịt sau khi ấp thành phôi dở chừng đem luộc lên thì sẽ như thế.''

Towa run rẩy chỉ tay vào thứ sinh vật trong bát của mình: ``A-Anh không cảm thấy sợ sệt hay gớm giếc gì sao!?''

Kuon nhíu mày nhìn nàng Ác quỷ: ``Ác ma như em mà sao nhát gan thế không biết\dots''

``K-Không có! Em là Ác ma nhưng em có trái tim nha! Anh không cảm thấy tiếc thương hay tội nghiệp cho chú vịt con sao!?''

Flare thở dài lắc đầu, trêu chọc hậu bối: ``\dots Chả trách người ta cứ gọi em là TMT (Towa-sama Maji Tenshi - Towa thật sự là thiên thần), Towa-sama ạ.''

``\textbf{Em không phải Thiên thần!}'' Towa ấm ức kêu lên.

Lúc này, Ryuuhou lanh chanh bóc phăng quả trứng vịt lộn trên tay mình, dùng thìa xúc trọn phần lòng đỏ béo ngậy cho vào miệng nhai nhòm nhèm đầy khoái chí, vừa ăn vừa chêm vào: ``Thương xót gì đâu! Món này là `thần dược' bổ dưỡng số một ở quê em đấy! Ăn vào là sung sức cả ngày luôn!''

Suzuno đứng bên cạnh nhẹ nhàng bóc một quả khác, khéo léo múc một chút phần cút dẻo đưa cho Aki, dịu dàng giải thích thêm: ``Aki-san thử một chút đi ạ, không đáng sợ đâu. Ăn cái này là phải húp phần nước ngọt đậm đà bên trong quả trứng trước, sau đó chấm phần nhân với chút muối tiêu tắc (quất) cay nồng, ăn kèm vài lá rau răm thơm phức thì mới chuẩn vị.''

Kuon gật gù tiếp lời em gái: ``Chuẩn đấy. Ayame-chan chuẩn bị trứng thế này là thiếu mất rau răm với muối tiêu tắc rồi, ăn thế xém ngon đi một nửa đấy.''

Aki, Flare và Towa nhìn ba anh em nhà Hikawa hào hứng chuyền tay nhau ăn những quả trứng vịt lộn một cách ngon lành mà toàn thân run lẩy bẩy trước sự ``máu lạnh'' và khẩu vị ``khủng'' của họ.

Ayame đứng bên cạnh cũng đổ mồ hôi hột, nụ cười gian xảo ban đầu biến mất hoàn toàn: ``Đ-Đây là nước đi mà em thật sự không ngờ tới đấy, Kuon-senpai à\dots\ Cả Ryuu-chan với Suzuno-chan nữa\dots''

Và thế là, cuộc thi dọa người của Ojou-san chính thức thất bại thảm hại, khi toàn bộ số trứng vịt lộn trên đĩa đều bị ba anh em nhà Hikawa và một Aki ``bị mua chuộc'' xử lý sạch bách!

\newpage

\subsection*{\phantomsection\label{ep-C9.IX}}
\addcontentsline{toc}{subsection}{{\sen{-Histoire 9-}} {\caviar Chả rươi}}
\stpt{-Histoire 9-}
\stcap{Chả rươi}
\vspace{-4pt}
\begin{center}
  {\caviar\fontsize{12pt}{12pt}\selectfont Nhân vật: \emp{kuon2} \emp{kaigou2} \emp{flare1} \emp{towa}}
\end{center}

``\textbf{Áaaaaaaaaaa!!! Cứu tôi với!!}''

Một tiếng la hét thất thanh vang vọng ra từ phòng bếp văn phòng.

Kuon đang ngồi đọc sách ngoài phòng khách vội vàng buông quyển sách xuống, chạy xông vào bếp. Đập vào mắt cậu là cảnh tượng Flare và Towa đang khúm núm ôm chặt lấy nhau ở góc tường, mặt mày cắt không còn một giọt máu.

Kuon lo lắng hỏi dồn: ``Có chuyện gì thế hai đứa!?''

Nàng Yêu tinh vàng run rẩy giơ ngón tay chỉ về chiếc rổ mây đặt trên bàn bếp: ``C-Có cái sinh vật gì ghê rợn lắm đang ngoe nguẩy trên bàn kìa anh!''

Nàng Ác ma mếu máo tiếp lời: ``Đáng sợ quá đi mất! Nhìn nó bầy nhầy lúc nhúc như hàng ngàn con sâu ấy!''

Flare nheo mắt nhìn kỹ lại rồi rùng mình: ``Không phải sâu\dots\ Chị thấy nó giống giun đất hơn á!''

Kuon tiến lại gần chiếc rổ. Bên trong là hàng trăm con sinh vật dạng giun màu hồng nhạt pha xanh lá đang bò lúc nhúc, ngoe nguẩy rộn ràng dưới lớp nước mỏng.

Đúng lúc này, Kaigou xắn tay áo bước ra từ phòng kho phía sau, trên tay bưng chiếc thau sạch cùng một đĩa vỏ quýt tươi và nắm lá thì là xanh mướt. Thấy hai cô gái \hll\ đang nép chặt vào nhau run bần bật, cô nàng Tổng quản \hls\ bật cười khẩy: ``Gì mà làm ầm ĩ lên thế hai cô nương? Mới thấy mấy con rươi bé tí này mà đã sợ chết khiếp rồi à?''

Kuon thở phào một hơi, quay sang nhìn em gái song sinh: ``Kaigou, rổ rươi này tươi rói đấy chứ. Em lấy từ khu chợ người Etsunan sáng nay đúng không?''

Kaigou đặt thau nước xuống bàn bếp cái *CỘC*, dõng dạc đáp: ``Chứ sao nữa Aniki. Em phải nhờ người quen cọc trước từ sáng sớm mới gom được một cân rươi béo múp míp thế này đấy. Thế mà hai đứa này vừa nhìn thấy đã hét toáng lên như thể phòng bếp bị quái vật ngoài hành tinh xâm lăng vậy.''

Flare tròn mắt: ``Rươi\dots? Chúng là loài vật gì vậy, Kuon-senpai, Kaigou-senpai?''

Kuon thong thả giải thích: ``Nó là một loài động vật thân mềm sống dưới tầng bùn đất ở các vùng nước lợ ven biển Etsunan.''

Towa nuốt nước bọt: ``G-Giun sống dưới bùn nước á!? Kinh dị quá!''

Kaigou chẳng chút sợ sệt, cô nhanh nhẹn trút rươi vào thau nước ấm rồi dùng tay khuấy nhẹ để rửa sạch bùn và rụng bớt lông tơ. Nhìn đôi tay tháo vát xử lý đống sinh vật nhung nhúc một cách tỉnh queo, cả Flare lẫn Towa đều toát mồ hôi hột vì thán phục độ dạn dĩ của cô.

Cô vừa thoăn thoắt làm vừa lên tiếng: ``Trông hình dạng bò ngoe nguẩy thế này có phần hơi dọa người thật, nhưng đây là `lộc trời' cực phẩm đấy. Rươi chỉ xuất hiện vào vài ngày ngắn ngủi vào dịp cuối thu đầu đông thôi, không phải lúc nào có tiền cũng mua được đâu.''

Kuon chêm vào: ``Đúng đấy. Một cân rươi tươi thế này có giá từ 2000 đến 3000 yên lận đấy. Lúc hiếm hàng hay trái mùa, giá của nó có thể vọt lên tận 5000 yên một cân cơ.''

Towa thốt lên kinh ngạc: ``Hả!? Sao đắt quái quỷ vậy!?''

Flare tò mò hỏi tiếp: ``Vậy hai người mua đống này về để làm gì thế ạ?''

Kaigou nhếch môi cười đắc ý, chỉ tay vào thau rươi:``Để làm chả rươi ăn chứ sao nữa.''

Towa hoảng hốt trợn tròn mắt, hai tay ôm chặt lấy đầu: ``Ý hai người là đem cái đống nhầy nhụa đang bò lúc nhúc đó đi làm thức ăn á!?''

Kuon gật đầu chắc nịch: ``Ừ. Đảm bảo với mấy đứa là ngon đến mức `quên lối về' luôn!''

Flare và Towa đồng thanh xua tay lia lịa, mặt xanh lét: ``K-Không đời nào bọn em ăn cái món sâu bọ kinh dị đó đâu!!''

Kaigou khoanh tay cười khẩy, giọng đầy thách thức: ``Cứ mạnh mồm đi! Chút nữa rán xong thơm nức mũi cả văn phòng rồi thì đừng có mà nuốt nước miếng xin xỏ đấy nhé!''

Nói là làm, cô bắt tay vào trổ tài. Với sức mạnh thể chất và sự khéo léo của mình, cô dùng đũa đánh nhuyễn rươi cùng thịt lợn xay, đập thêm vài quả trứng gà tươi, rắc hạt tiêu cay nồng, rồi trộn đều với hành hoa, thì là và linh hồn của món ăn: vỏ quýt thái sợi chỉ mỏng dính.

Chẳng mấy chốc, tiếng mỡ xèo xèo vang lên rộn rã trong chảo gang. Từng miếng chả rươi được Kaigou khéo léo thả vào dầu sôi, phồng to và chuyển dần sang màu vàng ruộm óng ả. Mùi thơm nức mũi của vỏ quýt nướng quyện với vị béo ngậy của thịt băm, trứng chiên và thảo mộc lan tỏa ngào ngạt khắp gian bếp, đánh bật hoàn toàn mọi mùi tanh tưởi ban đầu.

Trên bàn ăn, đĩa chả rươi nóng hổi, bốc khói nghi ngút, giòn rụm bên ngoài và mềm mọng bên trong đã được dọn ra thơm lừng.

Flare và Towa đứng cạnh bên không kìm được mà nuốt nước bọt cái *ỰC*.

Và kết quả là\dots

*SẬT\dots\ RỐP!*

Nàng Yêu tinh vàng nhắm mắt cắn thử một miếng chả rươi chấm chút nước mắm quất cay, rồi đôi mắt lập tức mở to sáng rực như đèn pha: ``\dots Trời đất ơi! Nó ngon đến phát điên luôn ấy!''

Nàng Ác ma bên cạnh cũng đang cuống quýt gắp miếng thứ ba vào bát: ``Không thể tin được luôn! Cái con vật trông kinh dị ban nãy làm sao lại có thể biến thành một món ăn đỉnh cao thế này cơ chứchứ\dots!''

Kaigou tựa lưng vào mép bàn, khoanh tay nhìn hai cô gái vừa ăn vừa xuýt xoa khen ngợi mà nhếch mép cười đắc thắng: ``Thấy chưa? Ban nãy ai bảo sống chết không thèm đụng đũa vào hả? Chị nói rồi, đồ ăn xứ Kawachi không thể nhìn cái mã bề ngoài mà vội phán xét đâu.''

Kuon cũng cười sảng khoái tiếp lời: ``Đã nói rồi mà còn không chịu tin cơ.''

Flare tấm tắc khen ngợi nghệ thuật nấu nướng của hai chị em: ``Em thật sự nể tài nấu nướng của Kaigou-senpai đấy! Kết hợp rươi với trứng, thịt băm và đặc biệt là vỏ quýt thái sợi\dots\ Mùi vỏ quýt thơm nồng bốc lên làm món ăn thanh tao hẳn, không còn chút mùi tanh nào luôn!''

Towa gật gù liên tục, mồm vẫn nhai nhồm nhoàm: ``Đúng đó! Vỏ ngoài chả giòn rụm, bên trong thì béo ngậy, hạt rươi cắn vào sần sật vui miệng dã man!''

Kuon vỗ tay cười lớn: ``Tuyệt cú mèo đúng không? Món này là một trong những niềm tự hào ẩm thực lớn nhất của người dân Kawachi bọn anh đấy!''

Kaigou thong thả rót ra bàn cho mỗi người một tách trà ấm, cất giọng sang sảng đầy hào sảng: ``Ăn uống là phải bạo gan mở lòng trải nghiệm như thế mới sướng chứ. Ở đời 'tốt gỗ hơn tốt nước sơn', vẻ ngoài kỳ lạ bao nhiêu thì hương vị bên trong lại tinh túy bấy nhiêu. Ăn cho đẫy vào rồi tí nữa nhớ dọn dẹp rửa bát phụ tôi đấy nhé hai cô nương!''

Cả Flare và Towa đều gật đầu lia lịa đồng tình, miệng vẫn không ngừng xuýt xoa. Quả thực, ẩm thực xứ sở của anh em nhà Hikawa chưa bao giờ ngừng làm cho các cô gái \hll\ đi từ sợ hãi tột cùng cho đến bất ngờ và mê mẩn không lối thoát!

\end{document}